\section{测度的逆极限}

在本节中，I 表示一个非空集合，它赋有预序关系，记作 \(i \leqslant j\)，并且关于此关系是有向的。回忆（GT, I, §4, No.~4），由 I 索引的拓扑空间逆系统是一个族 \((\mathrm{T}_i, p_{ij})\)，其中 \(\mathrm{T}_i\) 是拓扑空间，且 \(p_{ij}\) 是从 \(\mathrm{T}_j\) 到 \(\mathrm{T}_i\) 的连续映射，当 \(i \leqslant j\) 时；\(p_{ii}\) 是 \(\mathrm{T}_i\) 的恒等映射；并且 \(p_{ik} = p_{ij} \circ p_{jk}\) 当 \(i \leqslant j \leqslant k\) 时。设 T 是拓扑空间；称族 \((p_i)_{i \in I}\) 为 T 上的相容族，其中映射 \(p_i: \mathrm{T} \to \mathrm{T}_i\) 定义在 T 上，且族 \((p_i)_{i \in I}\) 满足 \(p_i = p_{ij} \circ p_j\) 对于 \(i \leqslant j\)。若 T 中不同的 \(x, y\) 存在 \(i \in I\) 使 \(p_i(x) \neq p_i(y)\)，则映射 \(\mathrm{T} = \varprojlim \mathrm{T}_i\) 将它们分离。拓扑空间 \(p_i\) 称为该逆系统的逆极限；其投影 \(\mathrm{T}_i\) 映射到 \((p_i)_{i \in I}\)，并且这一族投影相容。

